
Since $E(\cdot)$ is a continuous function of robot models, we assume the transition dynamics of the robot $E(\alpha)$ is differentiable and locally $L_1$-Lipschitz w.r.t. $\alpha$ in the sense that
\begin{equation}\label{eq:alpha:lipschitz:assumption}
\begin{split}
     \exists \varepsilon > 0,
    ||E(\alpha)(s,a) - E(\alpha^\prime)(s,a)|| \le L_1 |\alpha - \alpha^\prime|, 
    \forall s\in\mathcal{S}, \forall a\in\mathcal{A}, \forall |\alpha^\prime - \alpha | < \varepsilon    
\end{split}
\end{equation}
Moreover, we follow \cite{slbo} and assume the value functions of the robot models are $L_2$-Lipschitz w.r.t to some norm $||\cdot||$ in state space in the sense that
\begin{equation}\label{eq:state:lipschitz:assumption}
    |V^{\pi, M}(s) - V^{\pi, M}(s^\prime)| \le L_2 ||s-s^\prime||, \forall s, s^\prime \in \mathcal{S}
\end{equation}


By the assumption in Equation \eqref{eq:state:lipschitz:assumption} that the value functions of the robots are $L_2$-Lipschitz, as proven in Lemma 4.1 of \cite{slbo}, $\forall \phi > 0$, $M=E(\alpha), M^\prime=E(\alpha+\phi)$, we have 
\begin{equation}\label{eq:proof2}
\begin{split}
\exists \varepsilon > 0, \text{s.t. }
\mathop{\mathbb{E}}_{s} [ |V^{\pi,M}(s) - V^{\pi,M^\prime}(s)| ] 
& \le \frac{\gamma}{1-\gamma} L_2 \mathop{\mathbb{E}}_{(s,a)\sim \pi} || M(s,a) - M^\prime(s,a) || \\
& = \frac{\gamma}{1-\gamma} L_2 \mathop{\mathbb{E}}_{(s,a)\sim \pi} || E(\alpha)(s,a) - E(\alpha^\prime)(s,a) || \\
& \le \frac{\gamma}{1-\gamma} L_2 L_1 |\alpha - \alpha^\prime|, \forall |\alpha^\prime - \alpha | < \varepsilon 
\end{split}
\end{equation}
Since the value functions are differentiable by assumption in Equation \eqref{eq:alpha:lipschitz:assumption}, suppose $D = | \frac{\partial{V^{\pi,E(\alpha)}}}{\alpha} | $ is the absolute value of the derivative of $V^{\pi,E(\alpha)}$ w.r.t. to $\alpha$. 
According to Equation \eqref{eq:proof2}, $D$ is bounded.
We make a (strong) assumption that for all $\varphi\in [\alpha, \alpha+\xi]$, policy $\pi^*_{E(\varphi)}$ only achieves the best expected reward on robot $E(\varphi)$, so that
\begin{equation}
    \forall \varphi, \beta \in [\alpha, \alpha+\xi], 
    V^{\pi^*_{E(\varphi)},E(\varphi)} -
    V^{\pi^*_{E(\varphi)},E(\beta)} = D \cdot |\varphi - \beta| + o(|\beta-\varphi|^2)
    = D \cdot |\varphi - \beta| + o(\xi^2)
\end{equation}
From the definition, the value function of a policy $\pi$ on uniformly sampled robots $E(\beta)$ from $\beta\sim U(\alpha, \alpha+\xi)$ is
\begin{equation}
\begin{split}\label{eq:proof1}
    & \mathop{\mathbb{E}}_{ 
    \substack{
    \beta \sim U(\alpha, \alpha+\xi)
    }}
    \mathop{\mathbb{E}}_{ 
    \substack{(s_t, a_t) \sim \pi, E(\beta)
    }}
    \sum_{t} \gamma^t r_t \cdot\exp(h\beta) \\
    = & \mathop{\mathbb{E}}_{ 
    \substack{
    \beta \sim U(\alpha, \alpha+\xi)
    }}
    \exp(h\beta)
    \mathop{\mathbb{E}}_{ 
    \substack{(s_t, a_t) \sim \pi, E(\beta)
    }}
    \sum_{t} \gamma^t r_t  \\ 
    = & \mathop{\mathbb{E}}_{ 
    \substack{
    \beta \sim U(\alpha, \alpha+\xi)
    }}
    \exp(h\beta)
    V^{\pi,E(\beta)}  \\
    = & \frac{1}{\xi} \int_{\beta=\alpha}^{\alpha+\xi} \exp(h\beta)
    V^{\pi,E(\beta)} \dd{\beta}
\end{split}
\end{equation}
Supposed the $\pi$ that optimizes Equation \eqref{eq:proof1} is the policy that directly optimizes on one robot $E(\varphi), \varphi \in [\alpha, \alpha+\xi]$, i.e. $\pi^*_{E(\varphi)}$ optimizes Equation \eqref{eq:proof1}, then we treat $\pi^*_{E(\varphi)}$ and $V^{\pi^*_{E(\varphi)},E(\varphi)}$ as constants and the following variation should be zero
\begin{equation}
\begin{split}
    \delta \int_{\beta=\alpha}^{\alpha+\xi} \exp(h\beta)
    (V^{\pi^*_{E(\varphi)},E(\beta)} - V^{\pi^*_{E(\varphi)},E(\varphi)})  \dd{\beta} = 0
\end{split}
\end{equation}
which means
\begin{equation}
\begin{split}
    0 = & \int_{\beta=\alpha}^{\alpha+\xi} \frac{\partial}{\partial \beta} [ \exp(h\beta)
    (V^{\pi^*_{E(\varphi)},E(\beta)} - V^{\pi^*_{E(\varphi)},E(\varphi)}) ]  \dd{\beta}  \\ 
    \approx & \int_{\beta=\alpha}^{\varphi} \frac{\partial}{\partial \beta} [\exp(h\beta) D |\varphi-\beta| ] \dd{\beta} - \int_{\beta=\varphi}^{\alpha+\xi} \frac{\partial}{\partial \beta} [\exp(h\beta) D |\varphi-\beta| ] \dd{\beta} + o(\xi^2) \\
    = & D \int_{\beta=\alpha}^{\varphi} \frac{\partial}{\partial \beta} [\exp(h\beta) (\varphi - \beta) ] \dd{\beta} - D \int_{\beta=\varphi}^{\alpha+\xi} \frac{\partial}{\partial \beta} [\exp(h\beta) (\beta - \varphi) ] \dd{\beta} + o(\xi^2) \\
    = & \frac{e^{h\alpha}}{h} D [(h\varphi-1)(2e^{h(\varphi-\alpha)} - 1 - e^{h\xi}) + \\ & h(\alpha+\xi)e^{h\xi} - e^{h\xi} - h\varphi e^{h(\varphi-\alpha)} + e^{h(\varphi-\alpha)} - h\varphi e^{h(\varphi-\alpha)} + e^{h(\varphi-\alpha)} + h\alpha - 1 ] + o(\xi^2) \\
    = & \frac{e^{h\alpha}}{2h} D [(\alpha h^{2} \xi^{2} + 2 \alpha h \xi + 4 \alpha  + 2 h \xi^{2}  + 2 \xi) - (h^{2} \xi^{2} + 2 h \xi +4)\varphi ] +  o(\xi^2) 
\end{split}
\end{equation}
The last equation assumes $0 \le h(\varphi - \alpha) \le h\xi \ll 1$ and used second-order Taylor approximation of $e^y = 1 + y + \frac{1}{2}y^2 + o(y^2), y\in \mathbb{R}$. 
Then we have
\begin{equation}
\begin{split}
    \varphi & = \alpha + \frac{1}{2} \xi + \frac{\xi^2 h (2-\xi h)}{8 + 4\xi h + 2\xi^2 h^2} + o(\xi^2) \\
    & = \alpha + \frac{1}{2}\xi + \frac{1}{4} h\xi^2 + o(\xi^2)
\end{split}    
\end{equation}
Therefore, if
\begin{equation}
\begin{split}
    & \lim_{\xi \rightarrow 0}
    \left[
    \mathop{\arg\max}_{\pi}
    \mathop{\mathbb{E}}_{ 
    \substack{
    \beta \sim U(\alpha, \alpha + \xi) \\
    M_\beta = E(\beta)
    }}
    \mathop{\mathbb{E}}_{ 
    \substack{a_t \sim \pi(\cdot \mid s_t) \\ 
    s_{t+1} \sim M_\beta(\cdot \mid s_t, a_t)}}
    \sum_{t} \gamma^t r_t \exp(h \beta)
    \right] \\
    &
    = \pi_{M_\varphi}^* = \mathop{\arg\max}_{\pi}
    \mathop{\mathbb{E}}_{ 
    \substack{a_t \sim \pi(\cdot \mid s_t) \\ 
    s_{t+1} \sim M_\varphi(\cdot \mid s_t, a_t)}}
    \sum_{t} \gamma^t r_t,
\end{split}
\end{equation}
then when $\xi\rightarrow 0$, $\varphi = \alpha + \frac{1}{2}\xi + \frac{1}{4} h\xi^2 + o(\xi^2)$.
\qed